\documentclass[12pt,a4paper]{report}
\usepackage[left=2.5cm,right=4cm,top=2.5cm,bottom=2.5cm,marginparwidth=3cm]{geometry}
\usepackage{amsmath}
\usepackage{booktabs}
\usepackage{setspace}
\usepackage{xcolor}
\usepackage{todonotes}
\usepackage{microtype}
\usepackage[hidelinks]{hyperref}
\doublespacing

% Set this to the chapter number it will have in the full thesis.
\setcounter{chapter}{2}

% Drafting helpers. \wordtarget prints a reminder; \fixme adds a margin note.
\newcommand{\wordtarget}[1]{{\color{gray}\small\textit{[Target: about #1 words]}}\par}
\newcommand{\fixme}[1]{\todo[color=orange!30,size=\scriptsize]{#1}}
\newcommand{\draftstatus}{Draft 2, sent to supervisor 14 September 2026}

\begin{document}

\begin{center}
\fcolorbox{gray}{gray!8}{\parbox{0.9\linewidth}{\singlespacing\small
\textbf{Chapter draft.} \draftstatus.\\
Working title of thesis: \textit{Title of Your Thesis}. Author: Your Name.\\
Remove this box and the notes before merging into the full thesis.}}
\end{center}

\chapter{Evidence from Regional Archives}
\label{ch:archives}

\section*{Chapter overview}
\wordtarget{300}
Tell the reader what this chapter argues, how it connects to the previous chapter and what comes next.
\fixme{Link back to the research question stated in Chapter 1.}

\section{Introduction}
\wordtarget{800}
Frame the specific question the chapter answers.

\section{Sources and approach}
\wordtarget{1500}
Describe the material used in this chapter only. Refer to the full methods chapter rather than repeating it.

\begin{table}[ht]
\centering
\caption{Archival collections consulted (example).}
\label{tab:collections}
\begin{tabular}{lrl}
\toprule
Collection & Items read & Date range \\
\midrule
County minute books & 214 & 1880 to 1910 \\
Parish correspondence & 96 & 1885 to 1902 \\
Newspaper cuttings & 340 & 1878 to 1915 \\
\bottomrule
\end{tabular}
\end{table}

\section{Analysis}
\wordtarget{3000}
\subsection{First theme}
Build the argument with evidence from Table~\ref{tab:collections}.
\fixme{Add two more quotations here.}
\subsection{Second theme}
Contrast or extend the first theme.

\section{Discussion}
\wordtarget{1200}
What does this chapter add to the thesis as a whole?

\section{Chapter summary}
\wordtarget{300}
Summarise the chapter in a few sentences and point forward to the next one.

\section*{Notes to self}
\begin{itemize}
  \item Check page numbers for all archive references.
  \item Ask supervisor whether Section 3.4 belongs in the next chapter.
\end{itemize}

\begin{thebibliography}{9}
\bibitem{c1} A.~Historian, \emph{Local Government in Transition}, Example Press, 2011.
\bibitem{c2} B.~Scholar, Reading minute books, \emph{Archive Studies Example} 17 (2016) 55--72.
\end{thebibliography}

\end{document}
